% !TEX root = ../main.tex
% =====================================================================
% §3 The Disconnection--Reconnection Benchmark — gdn2vessel benchmark-only D&B 稿
% 真源：STORY_FRAMEWORK.md §3 故事弧 + LEADERBOARD_MATRIX.md（severity 网格 + 防泄漏）
% 公式照实现核实：
%   src/benchmark/synth_breaks.py  — SEVERITY_GRID (size_max 6/8/10/12, nb_deco=100)、
%       creatis create_disconnections 协议、apply_breaks/GapRecord/BreakResult
%   src/benchmark/metrics.py       — epsilon_beta0:173 / betti_error:234 / cldice_metric:307 /
%       _is_gap_closed:640 / success_rate:679 / reid_rate:773 / _check_reid_at_gap:711
%   src/datasets/precompute_benchmark.py — npz 离线冻结 + manifest + test-split-only(RED LINE 1)
% 数字一律 \TODO 占位（禁臆造），唯一已确证可写实 = SR=1.0@dice≈0.16 诊断（代码根因明确）。
% 区分度 claim hedge 至批2；双盲：无自引/机构。
% =====================================================================
\section{The Disconnection--Reconnection Benchmark}
\label{sec:benchmark}

The benchmark is the central methodological contribution of this paper. It
consists of (i) a synthetic-breakpoint protocol that injects reproducible
gaps into ground-truth vessel masks over a four-level severity grid
(Section~\ref{sec:benchmark:protocol}); (ii) a leakage-control design that
keeps evaluation sets held out and prevents the protocol from feeding
ground-truth topology into any evaluated method
(Section~\ref{sec:benchmark:leakage}); (iii) a family of connectivity
metrics for reconnection, cross-checked against standard topological measures
(Section~\ref{sec:benchmark:metrics}); and (iv) a documented diagnosis of a
metric failure mode, namely that the gap-closure success rate can be
saturated by over-prediction (Section~\ref{sec:benchmark:srdiag}).

\subsection{Synthetic breakpoint protocol}
\label{sec:benchmark:protocol}

\paragraph{Breakpoint injection.} The core disconnection algorithm follows
the published plug-and-play breakpoint procedure of
\cite{TODO_creatis}, which we adopt as a protocol reference rather than as a
competing method. Given a binary ground-truth vessel mask, the procedure
first computes a chessboard distance transform restricted to the morphological
skeleton, so that each skeletal pixel is tagged with the local vessel
thickness. Skeletal pixels are then bucketed into thickness tiers, and only
the three finest (thinnest) tiers are retained, since thin vessels are the
most plausible places for a real segmenter to break. A breakpoint draw
selects a tier with probability proportional to an exponential law that
favours thinner vessels: with $U$ retained tiers indexed by $i$,
\begin{equation}
  P(\text{tier } i) \;=\; \frac{2^{\,U-(i+1)}}{2^{\,U}-1},
  \qquad i = 0,\dots,U-1 ,
  \label{eq:tierprob}
\end{equation}
so thinner tiers are exponentially more likely to be broken. For each of
$N_{\mathrm{deco}}$ draws a skeletal point is sampled from the chosen tier and
a soft erasure artefact is stamped at that location: a solid disk of radius
$\rho \sim \mathcal{N}(s_{\max} / t,\; 4)$ (clipped at $\rho \geq 0$, where $t$
is the tier's thickness value and $s_{\max}$ the maximum gap-size parameter)
is sparsely sampled, Gaussian-smoothed with $\sigma = 0.8$, and thresholded at
$0.4$ to produce the erased region. All erasures are accumulated and the mask
is binarised at $\geq 0.1$. The procedure is purely subtractive: a
zero-leakage assertion enforces that every foreground pixel of the broken mask
was foreground in the original ground truth, so no spurious vessels are ever
hallucinated by the protocol.

\paragraph{Per-gap records.} On top of the \cite{TODO_creatis} procedure, our
suite additionally extracts a per-gap record for every connected component of
the accumulated erasure mask. Each record stores the gap centre, an
approximate radius (half the bounding-box diagonal), the gap-size parameter,
the smoothing $\sigma$, and the identities of the two dominant ground-truth
vessel segments bordering the gap. The segment identities are obtained by
connected-component labelling of the \emph{original} ground-truth mask
(8-connectivity), and they are what makes same-vessel re-identification
measurable (Section~\ref{sec:benchmark:metrics}). This per-gap metadata is
an addition of this paper; the reference procedure produces only the broken
mask.

\paragraph{Severity grid.} We expose the maximum gap-size parameter
$s_{\max}$ as a difficulty knob and define a four-level severity grid, with
the number of draws fixed at $N_{\mathrm{deco}} = 100$ throughout:

\begin{table}[t]
  \centering
  \begin{tabular}{lccl}
    \toprule
    Severity & $s_{\max}$ & $N_{\mathrm{deco}}$ & Provenance \\
    \midrule
    Easy    & 6  & 100 & this paper's extension \\
    Medium  & 8  & 100 & aligned with the reference protocol's official setting \\
    Hard    & 10 & 100 & this paper's extension \\
    Extreme & 12 & 100 & this paper's extension \\
    \bottomrule
  \end{tabular}
  \caption{Four-level severity grid. Only the Medium level is anchored to the
  reference protocol's published parameters ($s_{\max}=8$,
  $N_{\mathrm{deco}}=100$); the Easy, Hard, and Extreme levels are this
  paper's extensions and are flagged as such in the released code.}
  \label{tab:severity}
\end{table}

The grid is the benchmark's principal source of discriminative pressure: it
yields a \emph{severity-response} axis along which a method's reconnection
quality can be tracked as breakpoints widen
(Section~\ref{sec:benchmark}; leaderboard in Section~\ref{sec:benchmark}).

\paragraph{Reproducibility.} Every draw is governed by a fixed seed. Each
severity level derives its seed deterministically from a base seed
($\mathrm{seed} = 42$), so that a given (dataset, severity, seed) triple
always yields the identical set of gaps. The broken masks, the per-pixel
vessel-segment map, the per-gap records, and the preprocessed source image
are precomputed offline and frozen as compressed archives, indexed by a
manifest mapping each (dataset, severity) pair to its archive, gap count, and
timestamp. Because the breaks are frozen, every evaluated method is scored on
exactly the same masks, which is what makes the leaderboard comparable.

\subsection{Leakage-control design}
\label{sec:benchmark:leakage}

The benchmark enforces two non-negotiable separations.

\paragraph{Held-out evaluation sets.} Breakpoints are injected only into the
test split of each dataset; training images are never broken or evaluated
(for example, the DRIVE evaluation images are disjoint from the DRIVE
training images). The precompute stage reads ground-truth masks from the test
split exclusively, so no training sample can enter the benchmarked set.

\paragraph{The protocol never supervises an evaluated method with
ground-truth topology.} The synthetic breakpoint protocol is applied only to
construct the \emph{input} that an evaluated method must reconnect; the
ground-truth topology is used solely as the adjudicator at scoring time and is
never fed back to any method as a target or as a feature. Concretely, the
gap-injection procedure, the per-gap segment identities, and all evaluated
methods are input-derived, and the protocol is purely subtractive (the
zero-leakage assertion above guarantees the broken mask is a subset of the
original mask, so the protocol cannot leak new structure). Using the
ground-truth topology as the judge of reconnection quality is permitted;
using it to supervise a method under evaluation is not, and the frozen,
offline precompute makes this separation auditable.

\subsection{A family of connectivity metrics}
\label{sec:benchmark:metrics}

We report reconnection quality along a primary topological axis, two
reconnection-specific axes, and two standard topological cross-checks. Unless
noted, all metrics are computed from a predicted binary mask and the
ground-truth mask only, so that every baseline---whether or not it has a
re-identification head---is scored on the same footing.

\paragraph{Primary reconnection axis: $\varepsilon_{\beta_0}$.} The primary
axis is a topological continuity error based on the zeroth Betti number
$\beta_0$ (the number of connected components):
\begin{equation}
  \varepsilon_{\beta_0}
  \;=\;
  \frac{\bigl|\,\beta_0^{\text{pred}} - \beta_0^{\text{gt}}\,\bigr|}
       {\max\!\bigl(\beta_0^{\text{gt}},\, 1\bigr)} .
  \label{eq:epsbeta0}
\end{equation}
A perfectly reconnected mask matches the ground-truth component count and
attains $\varepsilon_{\beta_0} = 0$; broken predictions inflate the component
count and raise the error. We adopt $\varepsilon_{\beta_0}$ as the primary
axis because it measures connectivity directly rather than overlap, and
because it is more sensitive to true connectivity than the gap-level success
rate discussed below (Section~\ref{sec:benchmark:srdiag}).

\paragraph{Gap-closure success rate (SR).} \emph{This is a novel metric
contributed by this paper}; it is not part of the reference protocol, which
reports overlap and surface-distance measures together with
$\varepsilon_{\beta_0}$ only. SR uses the per-gap records to ask, for each
injected gap, whether it has been closed:
\begin{equation}
  \mathrm{SR}
  \;=\;
  \frac{N_{\text{closed}}}{N_{\text{gaps}}}
  \;=\;
  \frac{1}{N_{\text{gaps}}}
  \sum_{g=1}^{N_{\text{gaps}}}
  \mathbf{1}\!\left[\,\textsc{closed}(g)\,\right],
  \label{eq:sr}
\end{equation}
where $N_{\text{gaps}}$ is the number of injected gaps and
$\textsc{closed}(g)$ is the gap-closure test defined in
Eq.~\eqref{eq:closed} below. $\mathrm{SR} \in [0,1]$, with $1.0$ indicating
every gap closed. We deliberately report SR with a documented failure mode
(Section~\ref{sec:benchmark:srdiag}) rather than as a clean ranking axis.

\paragraph{Same-vessel re-identification rate.} A reconnection can close a
gap yet bridge the \emph{wrong} pair of vessels. To penalise such mistakes we
report a same-root re-identification rate, adapting identity-preservation
logic from multiple-object tracking (the IDF1 family) to the gap setting,
following the use of identity-aware matching in open-curve vessel tracking
\cite{TODO_snaketracker}:
\begin{equation}
  \mathrm{reid\text{-}rate}
  \;=\;
  \frac{1}{N_{\text{gaps}}}
  \sum_{g=1}^{N_{\text{gaps}}}
  \mathbf{1}\!\left[\,
    \textsc{seg}_L(g)\in\mathcal{C}(g)
    \;\wedge\;
    \textsc{seg}_R(g)\in\mathcal{C}(g)
  \,\right],
  \label{eq:reid}
\end{equation}
where $\textsc{seg}_L(g)$ and $\textsc{seg}_R(g)$ are the two ground-truth
segment identities recorded for gap $g$, and $\mathcal{C}(g)$ is the set of
ground-truth segment identities found under the predicted foreground within a
disk around the gap centre. A gap counts as correctly re-identified only if
\emph{both} bordering segments are bridged (for an internal single-segment
gap, the one recorded segment must appear); bridging a different vessel scores
$0$, never a partial credit, mirroring the strict identity logic of MOT IDF1.
Because Eq.~\eqref{eq:reid} is computed from the predicted mask against the
ground-truth segment map, it is available for every baseline; a separate
head-driven variant applies only to methods that expose an explicit
re-identification head and is reported as an additional column without
contaminating the shared axes.

\paragraph{Standard topological cross-checks.} To guard against a
self-validating custom metric, every reconnection number is cross-checked
against two off-the-shelf, label-free topology measures. The first is clDice
\cite{TODO_cldice}, the centreline Dice, computed on hard binary masks as
\begin{equation}
  \mathrm{clDice}
  \;=\;
  \frac{2\,\mathrm{T_{prec}}\,\mathrm{T_{sens}}}
       {\mathrm{T_{prec}} + \mathrm{T_{sens}}},
  \quad
  \mathrm{T_{prec}} = \frac{|S_{\text{gt}}\cap V_{\text{pred}}|}{|S_{\text{gt}}|},
  \quad
  \mathrm{T_{sens}} = \frac{|S_{\text{pred}}\cap V_{\text{gt}}|}{|S_{\text{pred}}|},
  \label{eq:cldice}
\end{equation}
where $V_{(\cdot)}$ is a foreground mask and $S_{(\cdot)}$ its skeleton. The
second is the Betti-number error. Writing $\beta_0$ for connected components
and $\beta_1$ for independent loops, recovered from the 2D
Euler--Poincar\'e relation $\beta_1 = \beta_0 - \chi$ via the Euler number
$\chi$, we report
\begin{equation}
  \beta_0^{\text{err}} = \bigl|\beta_0^{\text{pred}}-\beta_0^{\text{gt}}\bigr|,
  \quad
  \beta_1^{\text{err}} = \bigl|\beta_1^{\text{pred}}-\beta_1^{\text{gt}}\bigr|,
  \quad
  \text{Betti err} = \beta_0^{\text{err}} + \beta_1^{\text{err}}.
  \label{eq:betti}
\end{equation}
clDice and the Betti error use only standard definitions and provide an
independent topological reading against which our reconnection axes can be
sanity-checked. A segmentation sanity axis (Dice, IoU, AUC, sensitivity,
specificity) is reported alongside but is not treated as a reconnection
measure.

\subsection{Diagnosis: the gap-closure success rate is saturable by
over-prediction}
\label{sec:benchmark:srdiag}

Reporting a metric family is only useful if its failure modes are
characterized. We document one such failure mode for the gap-closure success
rate of Eq.~\eqref{eq:sr}, and it is the reason we relegate SR to a contextual
rather than a primary role.

\paragraph{The gap-closure test.} A gap $g$ centred at $(c_y, c_x)$ with
recorded radius $r$ is judged closed if the predicted mask contains
\emph{any} foreground pixel inside a disk of radius $r + 3$ around the gap
centre:
\begin{equation}
  \textsc{closed}(g)
  \;=\;
  \mathbf{1}\!\left[\,
    \exists\,(y,x):\;
    (y-c_y)^2 + (x-c_x)^2 \leq (r+3)^2
    \;\wedge\;
    \mathrm{pred}(y,x) > 0
  \,\right].
  \label{eq:closed}
\end{equation}
The test asks only whether the centre of the erased zone is covered, with a
small tolerance for spatial offset; it does not ask whether the covering
foreground actually re-establishes a path between the two severed segments.

\paragraph{The failure mode.} Eq.~\eqref{eq:closed} can be satisfied without
any genuine reconnection. A degenerate model that floods the image with
foreground---a single large blob---covers the centre of every gap and is
therefore credited with closing all of them, driving SR to its maximum. We
observe this directly: in a smoke test, a poorly fitted model with a Dice of
only $\approx 0.16$ attains a success rate of $1.0$. The two numbers are
incompatible only if SR is read as a reconnection score; they are perfectly
consistent once Eq.~\eqref{eq:closed} is recognised as a coverage test that a
blob satisfies trivially. The root cause is structural rather than a tuning
artefact: the gap-closure test in Eq.~\eqref{eq:closed} checks coverage of the
gap centre, not the restoration of connectivity, so over-prediction is
rewarded.

\paragraph{Consequence for the benchmark.} This diagnosis dictates our
metric protocol. We use the topological continuity error
$\varepsilon_{\beta_0}$ of Eq.~\eqref{eq:epsbeta0} as the primary reconnection
axis: a flooding blob that saturates SR \emph{fuses} the entire mask into one
component, which drives $\beta_0^{\text{pred}}$ far from $\beta_0^{\text{gt}}$
and is penalised by $\varepsilon_{\beta_0}$ rather than rewarded. SR is
reported only alongside a specificity / over-prediction context (so a high SR
accompanied by a collapsed specificity is read as flooding, not
reconnection), and never as a standalone ranking axis. Characterising a
metric's exploitable failure mode in this way follows the tradition of
metric-failure-mode analysis in segmentation evaluation
\cite{TODO_pixelwise,TODO_touchstone}, and it is, with the synthetic-breakpoint
protocol and the connectivity-metric family, one of the methodological
contributions of this suite.
